\chapter{Literature Review}
\label{chap:literature_review}

\section{Search Strategy}
\label{sec:search_strategy}

The literature review follows a systematic approach using the \textbf{Scopus} database as the primary source. The search is structured around the following keyword clusters derived from the research questions:

\begin{table}[htbp]
  \centering
  \caption{Keyword clusters for systematic literature search}
  \label{tab:keywords}
  \begin{tabular}{@{}ll@{}}
    \toprule
    \textbf{Cluster} & \textbf{Keywords} \\
    \midrule
    Resource domain   & CAPEX portfolio, capital expenditure, project portfolio management \\
    Resource planning & resource allocation, skill-based planning, workforce planning, bottleneck \\
    AI / Forecasting  & demand forecasting, machine learning, time series, LSTM, XGBoost \\
    Symbolic AI       & knowledge graph, ontology, ESCO, semantic web, skill taxonomy \\
    Hybrid AI         & hybrid AI, neuro-symbolic, knowledge-augmented ML \\
    Optimization      & resource-constrained scheduling, RCPSP, integer programming \\
    \bottomrule
  \end{tabular}
\end{table}

Each retrieved paper is assessed for relevance using an Excel screening sheet with the following columns:
\begin{itemize}
  \item \textbf{Relevance}: Low / Medium / High
  \item \textbf{Finding}: Key insight relevant to this research
  \item \textbf{Category}: Which research question or theme the paper addresses
\end{itemize}

Only papers rated Medium or High are included in the synthesis below.

% ─────────────────────────────────────────────────────────────────────────────
\section{CAPEX Portfolio and Resource Management}
\label{sec:lit_capex}

% TODO: synthesize papers on PPM resource management, CAPEX-specific literature

% ─────────────────────────────────────────────────────────────────────────────
\section{Skill Ontologies and Knowledge Representation}
\label{sec:lit_ontologies}

The European Commission's ESCO ontology provides a multilingual, hierarchical taxonomy of over 13,900 skills and competences linked to occupations and qualifications \citep{le2014esco}. Its formal structure makes it suitable as a \textit{semantic backbone} for skill-demand modeling, enabling consistent cross-project and cross-organizational skill matching.

% TODO: expand with knowledge graph applications in HR/workforce planning

% ─────────────────────────────────────────────────────────────────────────────
\section{The AI-4-SME Framework}
\label{sec:lit_ai4sme}

The AI-4-SME Framework \citep{peter2025ai4sme}, developed by Peter, Laurenzi, and Hinkelmann (2025) at FHNW School of Business, provides a structured, three-phase methodology (Design, Build, Run) for identifying and implementing AI solutions in enterprise contexts. Relevant to this thesis are:

\begin{itemize}
  \item The \textbf{KIA/DIA distinction}: the framework differentiates Knowledge-Intensive Activities — requiring problem solving, analysis, creativity and decision making — from Data-Intensive Activities requiring data analysis and transaction processing. CAPEX resource planning involves both: demand forecasting is a KIA (contextual, expert-judgment-dependent), while allocation computation is a DIA (high-volume, optimization-driven).
  \item The \textbf{hybrid AI build path}: the framework explicitly models a hybrid path combining data-based ML (TensorFlow, PyTorch) with knowledge-based systems (Neo4j, GraphDB, Stardog, RDFox, metaphactory, TopBraid Composer). This directly maps to the architecture proposed in this thesis.
  \item The \textbf{six W-question problem definition method}: What / Who / Why / When / Where / How — used to structure the point-of-view formulation for AI task design.
  \item \textbf{Governance considerations}: integration of GDPR and EU AI Act compliance into the Run phase, which is relevant for enterprise deployment of the recommendation engine.
\end{itemize}

\section{Hybrid AI: Symbolic and Subsymbolic Integration}
\label{sec:lit_hybrid_ai}

Hybrid AI combines the strengths of \textit{symbolic AI} — explicit, interpretable knowledge representations such as ontologies and rules — with \textit{subsymbolic AI} — statistical learning from data, including deep learning and generative models \citep{garcez2022neurosymbolic}. Gartner highlights hybrid AI as an emerging trend in enterprise AI adoption, particularly in domains where domain knowledge must constrain or enrich data-driven models \citep{gartner2024hype}.

% TODO: survey neuro-symbolic approaches, knowledge-augmented forecasting

% ─────────────────────────────────────────────────────────────────────────────
\section{Resource Demand Forecasting}
\label{sec:lit_forecasting}

% TODO: ML-based forecasting in project contexts, temporal models

% ─────────────────────────────────────────────────────────────────────────────
\section{Resource Allocation and Optimization}
\label{sec:lit_optimization}

% TODO: RCPSP literature, skill-based extensions, make/buy/upskill models

% ─────────────────────────────────────────────────────────────────────────────
\section{Gap Analysis}
\label{sec:lit_gap}

The reviewed literature reveals the following gap: while resource-constrained project scheduling is well-studied theoretically, and hybrid AI is gaining traction in enterprise contexts, no existing work combines (1) ESCO-grounded skill ontologies, (2) data-driven temporal demand forecasting, and (3) make/buy/upskill allocation recommendations in a unified framework for CAPEX portfolio management. This thesis addresses that gap.
